% !TeX root = ../../Book3.tex
\section{复曲线积分}
\label{b3:sec:0201}

圆周上的积分已在\subsecref{b3:subsec:010204}用于恢复函数。一般曲线可以绕过孔洞，也可以反复经过同一点；积分必须记住行走方向与次数。下面先把参数曲线、积分估计和端点公式分开处理。

\subsection{分段光滑曲线}
\label{b3:subsec:020101}

设 \(\gamma:[a,b]\rightarrow\mathbb C\) 连续。若存在有限分点，使它在每个闭子区间上为 \(C^1\)，就称它为\term[fenduan-guanghua-quxian]{分段光滑曲线}[piecewise smooth curve]。允许某些区间上的导数为零，不要求曲线单射。曲线的长度定义为
\[
 L(\gamma)=\int_a^b|\gamma'(t)|\,\mathrm{d}t.
\]
反向曲线 \(\gamma^-(t)=\gamma(a+b-t)\) 的长度相同；首尾相接的曲线可以拼接，其长度相加。后文的闭曲线包括自交曲线，不能将“闭曲线”与“简单闭曲线”混为一谈。

\ProseParagraphBreak

圆周 \(a+r e^{\mathrm{i}t}\)、\(0\le t\le2\pi\)，规定正向为逆时针。把参数变为 \(2t\) 会把圆周走两遍；它与原曲线有相同的像，却有不同的积分。对非单射曲线，像集不能代替参数化。

\subsection{围道积分}
\label{b3:subsec:020102}

设 \(f\) 在曲线像上连续，沿 \(\gamma\) 的积分为
\[
 \int_\gamma f(z)\,\mathrm{d}z=\int_a^b f\bigl(\gamma(t)\bigr)\cdot\gamma'(t)\,\mathrm{d}t.
\]
右端按实部与虚部分别积分，有限个分点不影响数值。保持方向的分段 \(C^1\) 换参由实变量换元公式保持积分；反向换参改变符号。闭曲线上的积分也称\term[weidao-jifen]{围道积分}[contour integral]。
\begin{proposition}\label{b3:prop:02-ml}
设 \(f\) 在分段光滑曲线 \(\gamma\) 的像上连续。则
\[
 \left|\int_\gamma f(z)\,\mathrm{d}z\right|\le L(\gamma)\cdot\max_{z\in\gamma}|f(z)|.
\]
积分对被积函数线性，对曲线拼接可加，且 \(\int_{\gamma^-}f\,\mathrm{d}z=-\int_\gamma f\,\mathrm{d}z\)。
\end{proposition}
\begin{proof}
由实积分的三角不等式，积分模不超过 \(\int_a^b|f(\gamma(t))|\cdot|\gamma'(t)|\,\mathrm{d}t\)。取出连续函数在紧曲线像上的最大值即得估计。线性和拼接从有限区间积分的相应性质得到；反向时令 \(s=a+b-t\)，导数与积分上下限的变化给出负号。
\end{proof}
在半径 \(r\) 的正向圆周上，直接参数计算给出
\[
 \int_{|z-a|=r}(z-a)^n\,\mathrm{d}z=
 \begin{cases}2\pi\mathrm{i},&n=-1,\\0,&n\in\mathbb Z\setminus\{-1\}.\end{cases}
\]
这项计算不要求负幂在圆盘中心有定义。它将成为 Laurent 系数提取的依据。

\ProseParagraphBreak

复积分并不是把实积分中的变量符号换成 \(z\)。对任意分段光滑曲线，\(\int_\gamma\overline z\,\mathrm{d}z\) 一般不能写成端点的某个复函数之差。若 \(\gamma\) 是正向简单闭曲线，由 \(z=x+\mathrm{i}y\) 展开，
\[
 \overline z\,\mathrm{d}z=x\,\mathrm{d}x+y\,\mathrm{d}y+\mathrm{i}(x\,\mathrm{d}y-y\,\mathrm{d}x).
\]
前两项是 \(\tfrac12\mathrm{d}(x^2+y^2)\)，闭路上为零；后两项由平面 Green 公式等于内部面积的 \(2\mathrm{i}\) 倍。若暂不使用 Green 公式，对圆周直接计算已可看到面积贡献。被积函数不全纯时，内部积分无法用 Cauchy 消边法消去。

\ProseParagraphBreak

长度估计在曲线族随参数变化时也有用，但必须同时控制长度。例如圆 \(|z|=r\) 上的 \(z^{-2}\) 模为 \(r^{-2}\)，长度为 \(2\pi r\)，粗估计随 \(r\downarrow0\) 增大；实际积分却恒为零。这表明长度估计给上界，不自动反映不同相位的抵消。反之，若一段长为 \(\ell_n\) 的曲线上误差上界为 \(\varepsilon_n\)，仅 \(\varepsilon_n\rightarrow0\) 也未必使积分误差趋于0，还需 \(\ell_n\varepsilon_n\rightarrow0\)。

计算路径积分时，先检查被积式有没有单值原函数，常比直接代入参数更有效。例如沿任意从 \(0\) 到 \(1+\mathrm{i}\) 的分段光滑曲线，\(z e^{z^2}\) 的积分均为 \((e^{2\mathrm{i}}-1)/2\)。路径形状可以很复杂，但原函数 \(e^{z^2}/2\) 已把形状信息消去。相反，\(\overline z\) 的积分不能这样处理。在以 \(0,1,1+\mathrm{i},\mathrm{i}\) 为顶点的正向方形上，逐边参数化给积分 \(2\mathrm{i}\)。这个数来自区域的面积，说明仅有连续性与光滑性都不足以保证闭路积分为零。

\ProseParagraphBreak
参数方向也必须与被积式分开判断。把同一曲线反向走一遍，积分变号；把参数换成一个严格递增的分段光滑参数，积分不变。若参数曲线将某段往返走过，去程与回程相消。以后把区域切开时出现的两条切口边，正是利用这条相消规则。边界上的点值无需处处有导数；分段光滑条件只要求每段参数有连续导数，使积分能够逐段定义。尖角不产生额外项。

\subsection{原函数}
\label{b3:subsec:020103}

\begin{proposition}[原函数的端点公式]\label{b3:prop:02-primitive-endpoints}
设 \(U\) 为开集，\(F\) 在 \(U\) 上全纯，\(F'=f\)。若 \(\gamma\) 是 \(U\) 内由 \(p\) 到 \(q\) 的分段光滑曲线，则
\[
 \int_\gamma f(z)\,\mathrm{d}z=F(q)-F(p).
\]
\end{proposition}
\begin{proof}
全纯函数的导数连续已由\cref{b3:thm:01-holomorphic-analytic}得到。在每个光滑参数区间上，链式法则给出 \((F\circ\gamma)'=f(\gamma)\cdot\gamma'\)。分别对实部和虚部应用微积分基本定理，再把有限个区间的端点差相加，中间端点相消。
\end{proof}
例如 \(z^n\) 在 \(\mathbb C\) 上有原函数 \(z^{n+1}/(n+1)\)，其中 \(n\ge0\)。函数 \(1/z\) 在任何避开零点的小圆盘上都有对数原函数，却在整个 \(\mathbb C\setminus\{0\}\) 上没有原函数：单位圆上的积分为 \(2\pi\mathrm{i}\)，违反端点公式。局部存在与全局存在的差别，来自闭曲线能否收缩而不是函数是否局部光滑。

\ProseParagraphBreak

为判断路径独立性，可以先在一个具体区域计算闭路周期。例如在右半平面 \(\operatorname{Re}z>0\)，第一章的对数主值满足 \((\operatorname{Log}z)'=1/z\)，所以由 \(p\) 到 \(q\) 的每条路径都给 \(\operatorname{Log}q-\operatorname{Log}p\)。在去心平面，沿从1到 \(-1\) 的上半圆与下半圆积分分别为 \(\pi\mathrm{i}\) 和 \(-\pi\mathrm{i}\)；拼合得到一个单位圆周期。两个端点相同而积分不同，正是不存在单值原函数的具体表现。

\ProseParagraphBreak

若定义域不连通，原函数在每个连通分支上可以任意加一个常数，不存在跨分支的端点公式。连通分支在平面开集中仍然是开集；构造原函数时固定的起点应逐分支选择。后文说“原函数相差常数”时，常数的作用域也必须是一个区域，而不是任意开集。
